% Generado por scripts/70_tablas_figuras.R a partir de resultados-v1. No editar a mano.
\begin{table}[htbp]
\centering
\small
\caption{Tabla 6.6—Intervalos robustos a violaciones de tendencias paralelas}
\label{tab:tabla_6_6}
\begin{threeparttable}
\begin{tabular}{@{}lccccc@{}}
\toprule
 & (1) & (2) & (3) & (4) & (5) \\
 & Convencional & $\bar{M}$ = 0,5 & $\bar{M}$ = 1 & M = 0 & M = M obs. \\
\midrule
\multicolumn{6}{@{}l}{\textit{Panel A. Con covariables}} \\
e = 0 & [$-$231; 248] & [$-$365; 433] & [$-$565; 707] & [$-$604; 502] & [$-$780; 678] \\
Promedio e = 0--5 & [$-$570; 170] & [$-$1.289; 888] & [$-$2.279; 1.878] & [$-$2.422; 1.599] & [$-$4.318; 3.496] \\
  M observado = 233 ha &  &  &  &  &  \\
\addlinespace
\multicolumn{6}{@{}l}{\textit{Panel B. Sin covariables}} \\
e = 0 & [$-$160; 456] & [$-$355; 852] & [$-$808; 1.343] & [$-$909; 418] & [$-$1.138; 653] \\
Promedio e = 0--5 & [$-$843; 348] & [$-$2.438; 1.271] & [$-$4.238; 3.046] & [$-$4.712; 1.478] & [$-$6.904; 3.672] \\
  M observado = 274 ha &  &  &  &  &  \\
\bottomrule
\end{tabular}
\begin{tablenotes}[flushleft]
\footnotesize
\item \textit{Notas:} Intervalos de confianza al 95 \% (ha por año) según Rambachan y Roth (2023). Columnas (2) y (3): magnitudes relativas, la violación posterior puede ser hasta $\bar{M}$ veces la mayor violación previa. Columnas (4) y (5): suavidad, la pendiente de la tendencia diferencial puede cambiar hasta M ha por año; M = 0 extrapola linealmente la tendencia previa y M observado es la mayor curvatura previa a la adopción. Estudio de eventos con referencia universal en g $-$ 1. Intervalos de suavidad calculados con el método condicional híbrido (C-LF), porque el método FLCI no pudo ejecutarse por una incompatibilidad de software. El valor de ruptura es cero en todos los casos: el intervalo convencional ya incluye el cero.
\item \textit{Fuente:} Cálculos propios.
\end{tablenotes}
\end{threeparttable}
\end{table}
